\documentclass[11pt]{article}
\usepackage[T1]{fontenc}
\usepackage{lmodern,amsmath,amssymb,amsthm,geometry,hyperref}
\geometry{margin=0.95in}
\hypersetup{colorlinks=true,urlcolor=blue,linkcolor=blue}
\newtheorem{theorem}{Theorem}
\newtheorem{corollary}[theorem]{Corollary}
\newtheorem{proposition}[theorem]{Proposition}
\newtheorem{lemma}[theorem]{Lemma}
\newtheorem{remark}[theorem]{Remark}
\title{Weighted Lucas Sums and Differences:\\
Coefficient Classification and Four OEIS Sequences}
\author{Qichao Wang}
\date{October 2026}
\begin{document}
\maketitle
\begin{abstract}
We classify positive representations of \(L_i-qL_j\),
\(qL_i-L_j\), and \(qL_i+L_j\) for integer \(q\).
For \(q\ge2\), the first form has only finitely
many collisions, with an explicit uniform index bound.
The other two have infinitely many repeated values
precisely for three arithmetic coefficient families:
integral Lucas quotients, integral Fibonacci quotients,
and half-Lucas coefficients. Signed quotients of Lucas
gaps characterize all finite collisions, and
global uniqueness is decidable by finite exact arithmetic.
For \(q=5\), we also determine all repeated values,
the two complete pairwise intersections, and exact
counting formulas for OEIS A191980--A191983.
\end{abstract}
Let $F_0=0,F_1=1$ and $L_0=2,L_1=1$, with $u_{n+2}=u_{n+1}+u_n$.
Use $F_{-n}=(-1)^{n+1}F_n$ and $L_{-n}=(-1)^nL_n$. Put
\[
S_q=\{L_i-qL_j>0:i,j\ge1\},\quad
T_q=\{qL_i-L_j>0:i,j\ge1\},\quad
P_q=\{qL_i+L_j:i,j\ge1\}.
\]
The recurrence gives
\begin{equation}\label{eq:identities}
L_{m+4}-L_m=5F_{m+2},\qquad
F_aL_b=F_{a+b}+(-1)^bF_{a-b}.
\end{equation}

\section*{Infinite-collision constructions}
\begin{theorem}[Lucas--Fibonacci quotient coefficients]
\label{thm:ratios}
Let $m>h\ge1$ have the same parity, and suppose
\[
q=
\begin{cases}
L_m/L_h\in\mathbb Z,& m,h\text{ odd},\\
F_m/F_h\in\mathbb Z,& m,h\text{ even}.
\end{cases}
\]
Then $q\ge2$, and for every $n\ge1$,
\begin{equation}\label{eq:ratio-difference}
qL_{n+m-h}-L_n
=qL_{n+m+h}-L_{n+2m}>0.
\end{equation}
In particular, $T_q$ and $P_q$ each contain infinitely
many repeated positive values.
\end{theorem}
\begin{proof}
For $m$ odd, the Lucas product formula gives
\[
L_{n+2m}-L_n=L_mL_{n+m},\qquad
L_{n+m+h}-L_{n+m-h}=L_hL_{n+m}.
\]
For $m,h$ even, the Lucas--Fibonacci identity instead gives
\[
L_{n+2m}-L_n=5F_mF_{n+m},\quad
L_{n+m+h}-L_{n+m-h}=5F_hF_{n+m}.
\]
In either case dividing the two identities
gives
\[
L_{n+2m}-L_n
=q(L_{n+m+h}-L_{n+m-h}),
\]
which is \eqref{eq:ratio-difference}.
The right-hand side of the first expression
$qL_{n+m-h}-L_n$ is strictly positive and increases
with $n$ because $q\ge2$ and $m-h\ge1$.
The equivalent sum collision is
\[
qL_{n+m-h}+L_{n+2m}
=qL_{n+m+h}+L_n.
\]
\end{proof}




\begin{theorem}[Half-Lucas collision coefficients]
\label{thm:half-lucas}
For every integer $r\ge0$, let $M=6r+3$ and
\[
q=\frac{L_M}{2}\in\mathbb Z.
\]
Then $T_q$ and $P_q$ each have infinitely many
repeated positive values. More precisely, for
every $x\ge M$,
\begin{equation}\label{eq:half-lucas}
qL_x-L_{x-M+1}
=qL_{x+3}-L_{x+M+1}>0,
\end{equation}
and
\[
qL_x+L_{x+M+1}
=qL_{x+3}+L_{x-M+1}.
\]
The coefficients begin
$2,38,682,12238,219602,\ldots$.
\end{theorem}
\begin{proof}
The Lucas residues modulo $2$ have period $3$,
so $L_M$ is even when $3\mid M$.
As $M$ is odd, the product identity gives
\[
L_{x+M+1}-L_{x-M+1}=L_ML_{x+1}.
\]
The recurrence also gives
$L_{x+3}-L_x=2L_{x+1}$.
Substituting $2q=L_M$ proves
\eqref{eq:half-lucas} and the sum identity.
As $x\ge M\ge3$,
$qL_x-L_{x-M+1}>0$; these values
strictly increase with $x$.
\end{proof}




\section*{Norm separation and finite collision bounds}
\begin{lemma}[Algebraic-norm separation]\label{lem:norm-separation}
For integers \(d\ge1\) and \(q\ge1\), with
\(\alpha=\varphi\) and \(\beta=-\alpha^{-1}\),
\begin{equation}\label{eq:normbound}
|\alpha^d-q|\ge\frac1{q+1}.
\end{equation}
\end{lemma}
\begin{proof}
The conjugate algebraic integers \(\alpha^d-q\)
and \(\beta^d-q\) have nonzero integral product:
\(\alpha^d=F_d\alpha+F_{d-1}\) is irrational.
Since \(|\beta^d-q|\le q+1\), the bound follows.
\end{proof}

For completeness, this estimate also makes every bounded
value search finite. If \(0<L_i-qL_j\le X\), then
\(i>j\) and Binet's formula gives
\[
L_j\le (q+1)(X+q+1)+1,\qquad
L_i\le X+qL_j.
\]
The bounds for \(0<qL_i-L_j\le X\) follow by exchanging
the indices when \(j>i\), and by
\((q-1)L_i\le X\) when \(j\le i\) and \(q\ge2\).
For sums, \(qL_i+L_j\le X\) directly bounds both indices.
Thus for any integer coefficient, all distinct values at most
\(X\), their multiplicities (when finite), and their ranks
are computed by finite exact arithmetic; negative coefficients
reduce by exchanging sum and difference forms.

\begin{proposition}[No simultaneous-shift positive collisions]
\label{prop:no-shift}
For every integer $q\ge2$, there are no
indices $i<k$, $j<\ell$ such that
\[
L_{n+i}-qL_{n+j}=L_{n+k}-qL_{n+\ell}>0
\quad\text{for every }n\ge0.
\]
\end{proposition}
\begin{proof}
An equality for all shifts implies, by Binet's formula,
both its $\alpha$- and $\beta$-coefficient equalities.
Put $h=k-i\ge1$, $e=\ell-j\ge1$, and $s=i-j$.
Positivity for large $n$ forces $\alpha^s>q$,
so $s\ge2$, and
\[
\alpha^s(\alpha^h-1)=q(\alpha^e-1)
\]
forces $e>h$, hence $e\ge2$.
Conjugation gives
\[
q|\beta^e-1|=|\beta|^s|\beta^h-1|
\le\alpha^{-2}(1+\alpha^{-1})=\alpha^{-1}.
\]
But $q|\beta^e-1|\ge2(1-\alpha^{-2})
=2/\alpha>\alpha^{-1}$, a contradiction.
\end{proof}


\section*{The global finite-collision bound}
A finite threshold is sufficient to compute any term. The stronger
statement below provides a bound independent of the value.
\begin{theorem}[Uniform complete collision certificate]
\label{thm:global}
Fix an integer $q\ge2$. Let $t=t(q)$ be the least
positive index such that
\begin{equation}\label{eq:t-bound}
L_t>200(q+1)^4.
\end{equation}
If two distinct positive representations satisfy
\[
L_i-qL_j=L_k-qL_\ell>0,\qquad i<k,
\]
then
\begin{equation}\label{eq:collision-cutoff}
\boxed{k\le6t.}
\end{equation}
Consequently every $S_q$ with $q\ge2$ has
only finitely many multiply represented positive values.
Its complete collision list and global uniqueness
can be determined by checking the finitely many
pairs $1\le j<i\le6t$.
\end{theorem}
\begin{proof}
Write $\alpha=\varphi$, $\beta=-\alpha^{-1}$.
The equality forces $j<\ell$.
Moreover, positivity gives $k>\ell$.
Put $d=k-\ell$, $h=k-i$ and $g=k-j$.
We prove $k\le6t$ by contradiction.
The bounds
\[
\frac12\alpha^n\le L_n\le2\alpha^n
\quad(n\ge0),\qquad
L_k-L_i\ge L_{k-2}
\]
follow from Binet's formula and monotonicity.
Since
\[
L_{k-2}\le L_k-L_i
=q(L_\ell-L_j)<qL_\ell,
\]
we have $\alpha^d<4q\alpha^2<\alpha^t$;
thus $1\le d<t$. Suppose $k>6t$.
Then $\ell=k-d>5t$.

Write
\[
x=L_k-qL_\ell
 =\alpha^\ell(\alpha^d-q)+\beta^k-q\beta^\ell.
\]
The norm argument in \eqref{eq:normbound} gives
$|\alpha^d-q|\ge1/(q+1)$.
Since $|\beta^k-q\beta^\ell|\le q+1$
and $\alpha^\ell>2(q+1)^2$, the positivity
of $x$ forces $\alpha^d-q>0$ and
\[
x\ge\frac{\alpha^\ell}{2(q+1)}.
\]
But $x<L_i\le2\alpha^i$, so
$\alpha^{\ell-i}<4(q+1)<\alpha^t$,
which gives $1\le h<2t$.

Next, the original equality can be rearranged as
\[
qL_j=L_i-L_k+qL_\ell
=\alpha^k A+\beta^k A',
\]
where
\[
A=\alpha^{-h}-1+q\alpha^{-d},\qquad
A'=\beta^{-h}-1+q\beta^{-d}.
\]
Both are conjugate algebraic integers, and
\[
|A'|\le\alpha^h+1+q\alpha^d
<3\alpha^{2t}<\alpha^{3t}.
\]
If $A=0$, then $qL_j=|\beta^k A'|<1$,
an impossibility.
Otherwise $AA'$ is a nonzero integer, so
$|A|\ge\alpha^{-3t}$.
As $k>6t$ and $qL_j>0$, this forces $A>0$ and
\[
qL_j>\tfrac12\alpha^{k-3t}.
\]
Using $qL_j\le2q\alpha^j$ gives
$\alpha^{k-j}<4q\alpha^{3t}<\alpha^{4t}$,
hence $g<4t$.

Finally, express the collision as
\[
0=L_k-L_i-qL_\ell+qL_j
=\alpha^k E+\beta^k E',
\]
with the conjugate algebraic integers
\[
E=1-\alpha^{-h}-q\alpha^{-d}+q\alpha^{-g},
\quad
E'=1-\beta^{-h}-q\beta^{-d}+q\beta^{-g}.
\]
The bounds on $d,h,g$ imply
\[
|E'|\le1+\alpha^{2t}+q\alpha^t+q\alpha^{4t}
<\alpha^{6t}.
\]
If $E\ne0$, the integral norm $EE'\ne0$
gives $|E|\ge1/|E'|$.
The displayed equality now implies
$\alpha^{2k}\le|E'|^2<\alpha^{12t}$,
contrary to $k>6t$.

If $E=0$, the equality also forces $E'=0$.
It then holds after every simultaneous translation
of all four indices. Because $\alpha^d-q>0$,
the translated differences are positive for all
sufficiently large translations, contradicting
Proposition~\ref{prop:no-shift}.
Thus neither case is possible, and $k\le6t$.
\end{proof}




\section*{The stationary collision bound}
\begin{theorem}[Finite stationary pattern criterion]\label{thm:stationary}
Fix any integer $q\ge2$ and choose $t$ as in
\eqref{eq:t-bound}. If
\begin{equation}\label{eq:general-T-collision}
qL_k-L_\ell=qL_i-L_j>0,\qquad i<k,
\end{equation}
and $k>12t$, then all four indices lie
within $6t$ of $k$, apart from a possible
upper extension of at most $t$, and the same equality
holds after \emph{every} simultaneous nonnegative shift
of its indices.
Consequently, the set $T_q$ has infinitely many
multiply represented positive values if and only if
there are indices with
\[
K=12t+1,\quad K-3t\le i<K,\quad
K-t\le\ell\le K+t,\quad
K-6t\le j<\ell
\]
such that
\begin{equation}\label{eq:finite-shift-test}
0<qL_K-L_\ell=qL_i-L_j,\qquad
qL_{K+1}-L_{\ell+1}=qL_{i+1}-L_{j+1}.
\end{equation}
This is a finite, exact, necessary-and-sufficient
integer test for the existence of infinitely
many positive collisions in $T_q$.
Every collision not belonging to a stationary
family has larger first index $k\le12t$.
\end{theorem}
\begin{proof}
Equation \eqref{eq:general-T-collision} forces
$j<\ell$ and
$q(L_k-L_i)=L_\ell-L_j$.
The bounds $L_{n-2}\le L_n-L_{n-1}$
and $\frac12\alpha^n\le L_n\le2\alpha^n$ yield
\[
qL_{k-2}\le L_\ell,\qquad
L_{\ell-2}\le qL_k.
\]
Consequently, with $\delta=\ell-k$,
\[
\alpha^{-\delta}<4\alpha^2/q<\alpha^t,
\qquad
\alpha^\delta<4q\alpha^2<\alpha^t,
\]
so $|\delta|<t$.
Assume $k>12t$.
Set
\[
x=qL_k-L_\ell
=\alpha^k A+\beta^k A',
\quad
A=q-\alpha^\delta,\quad A'=q-\beta^\delta.
\]
The conjugate algebraic integers $A,A'$ have
nonzero integral norm, including when $\delta=0$
(since then $A=q-1$).
As $|A'|<\alpha^{2t}$,
we have $|A|>\alpha^{-2t}$.
Positivity of $x$ and $k>12t$ imply
$A>0$ and $x>\frac12\alpha^{k-2t}$.
Since $x<qL_i\le2q\alpha^i$,
the gap $h=k-i$ satisfies $0<h<3t$.

Next,
\[
L_j=qL_i-qL_k+L_\ell
=\alpha^k B+\beta^k B',
\]
where $B=q\alpha^{-h}-q+\alpha^\delta$
and $B'=q\beta^{-h}-q+\beta^\delta$.
The conjugates are algebraic integers and
$|B'|<\alpha^{5t}$. If $B=0$,
then $L_j<1$, a contradiction;
otherwise their integral norm gives
$|B|>\alpha^{-5t}$.
As $k>12t$ and $L_j>0$,
we obtain $L_j>\frac12\alpha^{k-5t}$.
Thus $g=k-j<6t$.
Since $j<\ell$, also $g>-t$.

Now the collision is
\[
0=qL_k-qL_i-L_\ell+L_j
=\alpha^k E+\beta^k E',
\]
where
\[
E=q-q\alpha^{-h}-\alpha^\delta+\alpha^{-g},
\quad
E'=q-q\beta^{-h}-\beta^\delta+\beta^{-g}.
\]
Both are algebraic integers and
$|E'|<\alpha^{8t}$.
If $E\ne0$, the integral norm gives
$|E|>1/|E'|$ and the equality forces
$\alpha^k\le |E'|<\alpha^{8t}$,
contradicting $k>12t$.
Therefore $E=E'=0$, and the identity
holds under all simultaneous shifts.
The positive leading coefficient $A$
makes its common values positive and unbounded
under increasing shifts.

Every infinite sequence of distinct collision
values must have unbounded $k$, and so contains
a stationary identity of the preceding form.
Translate its indices to $k=K=12t+1$.
The established gap bounds give precisely
the finite ranges in \eqref{eq:finite-shift-test}.
The leading coefficient $A>0$ ensures
positivity at $K$. Conversely, the two
equalities in \eqref{eq:finite-shift-test}
give consecutive zeros of a Lucas-recursive
difference. They therefore give the identity
at all shifts, and its positive values become
arbitrarily large. This proves the criterion
and the asserted bound for nonstationary collisions.
\end{proof}

\begin{proposition}[Finite certificates for all three forms]
\label{prop:finiteTP}
Fix \(q\ge2\) and \(t\) as in \eqref{eq:t-bound}.
If \(T_q\) has no stationary collision, every positive
collision of \(T_q\) or \(P_q\) has all indices at most
\(13t\). If a stationary collision exists, both forms
have a positive collision with all indices at most
\(13t+1\). In either case global uniqueness is decided
by a finite search. For \(S_q\), the bound is \(6t\).
For \(q=1\) all three forms have infinitely many
collisions; for \(q=0\), \(T_0\) is empty and \(S_0,P_0\)
have infinite representation multiplicities. For
\(q=-m<0\), \(T_q\) is empty, \(S_q=P_m\), and
\(P_q=S_m\) after exchanging index roles.
\end{proposition}
\begin{proof}
Theorem~\ref{thm:global} settles \(S_q\).
When \(T_q\) has no stationary collision,
Theorem~\ref{thm:stationary} puts its larger first index
at most \(12t\). Positivity gives \(L_j<qL_i\), and
hence \(j<i+t\), so all indices are at most \(13t\).
A sum collision with unequal first indices
has a corresponding signed difference
\[
v=qL_k-L_j=qL_i-L_\ell.
\]
If \(v>0\), this is a collision in \(T_q\); if
\(v<0\), its negative is one in \(S_q\).
The case \(v=0\) would give two different
representations of the same Lucas ratio \(q>1\),
contrary to Lemma~\ref{lem:lucas-ratio}.
Thus the same index bound holds for \(P_q\).
If \(T_q\) has a stationary collision,
Theorem~\ref{thm:stationary} supplies one at
\(k=12t+1\) with all other indices at most
\(13t+1\); swapping subtracted indices
gives the required sum collision.
The remaining signs follow from the recurrence
and the definitions.
\end{proof}

\section*{Explicit classification of infinite collisions}
The quotient and half-Lucas identities exhaust the infinite-collision
coefficients. We classify the index gaps of every stationary identity.

\begin{lemma}[Rigidity of the two index gaps]\label{lem:gap-rigidity}
Suppose \(q\ge2\), \(i<k\), \(j<\ell\), and
\begin{equation}\label{eq:stationary-difference}
q(L_{n+k}-L_{n+i})=L_{n+\ell}-L_{n+j}
\qquad(n\ge0).
\end{equation}
Set \(h=k-i\), \(e=\ell-j\), and \(d=i-j\).
The only possibilities are:
\begin{enumerate}
\item \(h=2r\), \(e=2m\), \(m>r\ge1\), with
 \(m\equiv r\pmod2\), \(d=m-r\), and either
 \(q=L_m/L_r\) for \(m,r\) odd or \(q=F_m/F_r\)
 for \(m,r\) even;
\item \(h=1\), \(e=2m\), \(m\ge3\) odd,
 \(d=m+1\), and \(q=L_m\);
\item \(h=3\), \(e=2m\), \(m\ge3\) odd and divisible
 by \(3\), \(d=m-1\), and \(q=L_m/2\);
\item \((h,e,d,q)=(1,3,2,2)\) or \((2,3,0,2)\).
\end{enumerate}
Every identity on this list is realized, with positive
values and arbitrarily large simultaneous shifts.
\end{lemma}
\begin{proof}
Let \(\alpha=\varphi\), \(\beta=-\alpha^{-1}\), and set
\[
C(s)=L_s-1-(-1)^s,\qquad
\Phi(s)=\frac{1+\alpha^{-s}}{1-\alpha^{-s}}\quad(s\ge1).
\]
The two conjugate Binet coefficient equations of
\eqref{eq:stationary-difference} give
\[
q\alpha^i(\alpha^h-1)=\alpha^j(\alpha^e-1),
\qquad
q\beta^i(\beta^h-1)=\beta^j(\beta^e-1).
\]
Their norms yield
\begin{equation}\label{eq:gap-norm}
q^2 C(h)=C(e),\qquad i\equiv j\pmod2.
\end{equation}
In particular \(e>h\), and \(d=i-j\) is even.
Dividing the two Binet equations, and using
\(\beta=-\alpha^{-1}\), gives the exact relation
\begin{equation}\label{eq:rigidity-ratio}
\alpha^{2d}=\alpha^{e-h}
\frac{\Phi(h)^{[h\text{ odd}]}}
     {\Phi(e)^{[e\text{ odd}]}}.
\end{equation}
The function \(\Phi(s)\) decreases strictly to \(1\);
direct calculation gives
\begin{equation}\label{eq:rigidity-special}
\Phi(1)=\alpha^3,\qquad \Phi(3)=\alpha.
\end{equation}

If \(h,e\) are both even, \eqref{eq:rigidity-ratio}
gives \(2d=e-h\). Write \(h=2r\), \(e=2m\).
Since \(d\) is even, \(m\equiv r\pmod2\).
For even indices,
\[
C(2s)=L_{2s}-2=
\begin{cases}
L_s^2,&s\text{ odd},\\
5F_s^2,&s\text{ even}.
\end{cases}
\]
Equation \eqref{eq:gap-norm} gives the two quotient
conditions in case~1.

If both gaps are odd, put \(u=(e-h)/2\).
Equation \eqref{eq:rigidity-ratio} becomes
\[
\alpha^{2(d-u)}=\Phi(h)/\Phi(e).
\]
When \(h\ge3\) this ratio lies strictly between
\(1\) and \(\Phi(3)=\alpha<\alpha^2\),
which is impossible for an integral exponent \(2(d-u)\).
When \(h=1\), the ratio is at least
\(\Phi(1)/\Phi(3)=\alpha^2\), with equality at \(e=3\);
for \(e>3\) it lies strictly between
\(\alpha^2\) and \(\alpha^3<\alpha^4\).
Therefore \(h=1,e=3\); then \(d=2\) and
\eqref{eq:gap-norm} gives \(q=2\).

If \(h\) is odd and \(e\) is even, the extra factor in
\eqref{eq:rigidity-ratio} is \(\Phi(h)\).
For odd \(h\ge5\), the right side lies strictly
between consecutive powers \(\alpha^{e-h}\) and
\(\alpha^{e-h+1}\), of which only the second has an
even exponent. This is impossible. For \(h=1\),
\eqref{eq:rigidity-special} gives \(2d=e+2\);
for \(h=3\) it gives \(2d=e-2\).
Putting \(e=2m\) in the norm equation then gives,
respectively, \(q=L_m\) or \(q=L_m/2\) when \(m\)
is odd. Indeed \(m\) even would require an irrational
positive multiple of \(\sqrt5\) to be the integer \(q\).
In the second case, integrality of \(L_m/2\) is
equivalent to \(3\mid m\); the condition \(q\ge2\)
also excludes \(m=1\).

Finally, for even \(h\), odd \(e\),
\eqref{eq:rigidity-ratio} has the factor \(1/\Phi(e)\).
For \(e\ge5\), the exponent lies strictly between
the consecutive integers \(e-h-1\) and \(e-h\);
the former is even, and no even exponent is possible.
Thus \(e=3\), necessarily \(h=2\), and
\eqref{eq:rigidity-special} gives \(d=0\) and \(q=2\).

For sufficiency, case~1 is the quotient identity
\eqref{eq:ratio-difference}. Cases~2 and~3 follow
from \(L_{n+2m}-L_n=L_mL_{n+m}\) for odd \(m\),
together with
\(L_{n+1}-L_n=L_{n-1}\) and
\(L_{n+3}-L_n=2L_{n+1}\).
The last two cases are the corresponding \(q=2\)
identities. All hold for every sufficiently large
common shift, with strictly increasing positive values.
\end{proof}

\begin{theorem}[Explicit infinite-collision coefficient classification]
\label{thm:explicit-infinite}
For every integer \(q\ge2\), the set \(T_q\) has
infinitely many multiply represented positive values
if and only if \(q\) belongs to the following union:
\begin{equation}\label{eq:explicit-coefficients}
\boxed{\begin{aligned}
&\left\{L_m/L_r\in\mathbb Z:
  m>r\ge1,\ m,r\text{ odd}\right\}\\
{}\cup{}&
\left\{F_m/F_r\in\mathbb Z:
  m>r\ge2,\ m,r\text{ even}\right\}\\
{}\cup{}&
\left\{L_{6r+3}/2:r\ge0\right\}.
\end{aligned}}
\end{equation}
The same equivalence holds for \(P_q\).
Every other integer \(q\ge2\) has only finitely many
collisions in \(T_q\) and \(P_q\).
\end{theorem}
\begin{proof}
Each coefficient in \eqref{eq:explicit-coefficients}
has an infinite family of positive collision values
by Theorems~\ref{thm:ratios} and~\ref{thm:half-lucas}.
Conversely, if \(T_q\) has infinitely many collision
values, Theorem~\ref{thm:stationary} forces a collision
holding under every common index shift.
Lemma~\ref{lem:gap-rigidity} puts its coefficient
in one of the three families: cases~2 and~4
are subsumed by the first or third family, and
case~3 is the third family.
The equivalence for \(P_q\) and finiteness outside
these families follow from
Proposition~\ref{prop:finiteTP}.
\end{proof}

\section*{Finite-collision coefficient spectra}
The preceding global bounds can also classify the coefficients
with finitely many collisions without computing any ordered
set of positive values. The sole arithmetic objects needed
are differences of two Lucas numbers and the sign of a
two-by-two determinant.

For integers \(1\le a<b\) put \(D(a,b)=L_b-L_a>0\).
Given two index pairs \(a<b\) and \(c<d\), define
\[
X=D(a,b),\quad Y=D(c,d),\qquad
\Delta(a,b;c,d)=L_aL_d-L_bL_c.
\]
The integrality condition \(Y\mid X\) generates
the candidate coefficient \(q=X/Y\).

\begin{lemma}[Uniqueness of positive Lucas ratios]
\label{lem:lucas-ratio}
For \(a>b\ge1\), \(c>d\ge1\), the equality
\(L_a/L_b=L_c/L_d\) implies \((a,b)=(c,d)\).
\end{lemma}
\begin{proof}
Set \(R_h(n)=L_{n+h}/L_n\), \(\alpha=\varphi\).
Binet's formula gives, for \(n\ge2,h\ge1\),
\begin{equation}\label{eq:ratio-interval}
\frac4{3\alpha}\alpha^h
\le R_h(n)\le \frac2\alpha\alpha^h,
\end{equation}
since
\[
\frac{1-\alpha^{-6}}{1+\alpha^{-4}}
=\frac4{3\alpha},\qquad
\frac{1+\alpha^{-6}}{1-\alpha^{-4}}
=\frac2\alpha.
\]
The intervals for distinct \(h\) are disjoint,
because \(4/3>2/\alpha\). For fixed \(h\),
\[
L_mL_{n+h}-L_nL_{m+h}
=5(-1)^{n+1}F_{m-n}F_h\ne0
\qquad(m>n\ge2).
\]
Thus all ratios with \(n\ge2\) are distinct.

For \(n=1\), write \(R_h(1)=L_s\),
\(s=h+1\ge2\). Since
\[
(1-\alpha^{-4})\alpha^s
\le L_s\le(1+\alpha^{-4})\alpha^s,
\]
the inequalities
\(\alpha(1-\alpha^{-4})>2/\alpha\)
and \(1+\alpha^{-4}<4/3\) show that
\(L_s=R_e(m)\), \(m\ge2\), would force \(e=s\).
But then the Lucas product formula gives
\[
L_sL_m-L_{m+s}=(-1)^sL_{m-s}\ne0,
\]
including \(m<s\) via the standard extension
to negative Lucas indices.
Ratios with both denominators \(L_1\)
are strictly increasing with their numerators.
\end{proof}

\begin{theorem}[Signed-gap coefficient spectra]
\label{thm:signed-gap}
Fix \(q\ge2\). The following statements are exact
necessary-and-sufficient conditions for the existence
of a multiply represented \emph{positive} value:
\begin{enumerate}
\item \(S_q\) has a collision if and only if some
 \(1\le a<b\), \(1\le c<d\) satisfy
 \(D(a,b)=qD(c,d)\) and \(\Delta(a,b;c,d)>0\).
\item \(T_q\) has a collision if and only if such
 indices satisfy \(\Delta(a,b;c,d)<0\).
\item \(P_q\) has a collision if and only if such
 indices exist, regardless of the sign of \(\Delta\).
\end{enumerate}
All three tests can be restricted to the common finite
index window
\begin{equation}\label{eq:signed-cutoff}
\boxed{1\le a,b,c,d\le 13t+1,}
\end{equation}
where \(t=t(q)\) is defined by \eqref{eq:t-bound}.
The two strict determinant signs construct all
collision witnesses using exact integer arithmetic.
\end{theorem}
\begin{proof}
If \(D(a,b)=qD(c,d)\), then the two unequal pairs
of indices give the sum collision
\[
qL_c+L_b=qL_d+L_a>0.
\]
Furthermore,
\[
L_a-qL_c=L_b-qL_d=\frac{\Delta}{D(c,d)},
\]
and
\[
qL_c-L_a=qL_d-L_b=-\frac{\Delta}{D(c,d)}.
\]
This proves that positive, negative, and zero signs
have precisely the asserted consequences.
Conversely, any collision in \(S_q\) has first
indices \(a<b\) and second indices \(c<d\);
its equality and positivity give the divisor equation
and \(\Delta>0\). For a \(T_q\) collision, label the
first indices \(c<d\) and second indices \(a<b\);
the same operations yield \(\Delta<0\).
For \(P_q\), the two first and second indices
necessarily run in opposite directions;
relabeling them gives the divisor equation.
Thus all three tests are equivalent, not merely
sufficient.

If \(\Delta=0\), the divisor equation implies
\(L_a=qL_c\) and \(L_b=qL_d\). These are two
distinct representations of the same Lucas ratio
\(q\ge2\), contrary to Lemma~\ref{lem:lucas-ratio}.
Thus no zero-determinant coefficient candidate exists.

For the common index bound, Theorem~\ref{thm:global}
places all \(S_q\) collision indices at most \(6t\).
If \(q\) has no stationary \(T_q\) collision,
Proposition~\ref{prop:finiteTP} places every
\(T_q\) and \(P_q\) collision inside \(13t\).
Otherwise Theorem~\ref{thm:stationary} supplies
a positive stationary \(T_q\) collision at
\(k=12t+1\), whose other three indices belong to
\([1,13t+1]\); exchanging subtracted indices
also supplies a \(P_q\) collision inside that window.
In all cases, if a collision exists, a certificate
satisfying \eqref{eq:signed-cutoff} exists.
\end{proof}

Define two explicitly parametrized sets of
positive integer coefficients:
\begin{equation}\label{eq:sign-spectra}
\boxed{\quad
\mathcal C_{\pm}
=\left\{\frac{D(a,b)}{D(c,d)}\in\mathbb Z_{\ge2}:
\, a<b,\ c<d,\ \pm\Delta(a,b;c,d)>0\right\}.
\quad}
\end{equation}
All four indices are positive. Thus \(S_q\) has a
positive collision precisely for \(q\in\mathcal C_+\);
\(T_q\) precisely for \(q\in\mathcal C_-\);
and \(P_q\) precisely for
\(q\in\mathcal C_+\cup\mathcal C_-\).
Let \(\mathcal I\) be the three explicit coefficient
families in \eqref{eq:explicit-coefficients}.
By Theorem~\ref{thm:explicit-infinite}, the
\emph{finitely collided but nonunique} coefficients are
\[
\begin{array}{c|l}
S & \mathcal C_+\\
T & \mathcal C_-\setminus\mathcal I\\
P &(\mathcal C_+\cup\mathcal C_-)
       \setminus\mathcal I.
\end{array}
\]
The complements give global uniqueness, while
\(\mathcal I\) gives infinitely many collision values
for \(T\) and \(P\). For \(S\), all collisions are
finite at every \(q\ge2\).

\begin{corollary}[Sum uniqueness from both differences]
\label{cor:sum-both}
For every integer \(q\ge2\), all positive values
in \(P_q\) have unique index-pair representations
if and only if this holds simultaneously in
\(S_q\) and \(T_q\).
\end{corollary}
\begin{proof}
Theorem~\ref{thm:signed-gap} partitions all positive
\(P_q\) collision witnesses into the two strict
signs, respectively equivalent to \(S_q\) and \(T_q\)
collisions. The zero sign is impossible by
Lemma~\ref{lem:lucas-ratio}.
\end{proof}

\begin{proposition}[Equal-gap parity rule]\label{prop:equal-gap}
For \(h\ge1\) and \(1\le c<a\), suppose
\[
q=\frac{L_{a+h}-L_a}{L_{c+h}-L_c}
\in\mathbb Z_{\ge2}.
\]
Then \(P_q\) has repeated positive values.
If \(c\) is odd, \(S_q\) has a positive collision;
if \(c\) is even, \(T_q\) has a positive collision.
In particular, each of the infinite families
\[
\boxed{\quad
q=\frac{L_{3r}}2\ (r\ge1),\qquad
q=\frac{L_{4r+2}}3\ (r\ge1),\qquad
q=\frac{F_{3r}}2\ (r\ge2)
\quad}
\]
consists entirely of coefficients for which
\(S_q\) has a finite nonempty collision set.
\end{proposition}
\begin{proof}
The Lucas--Fibonacci product identity gives
\begin{equation}\label{eq:gap-determinant}
\boxed{\quad
L_aL_{c+h}-L_cL_{a+h}
=5(-1)^{c+1}F_{a-c}F_h.\quad}
\end{equation}
For \(a>c\) and \(h\ge1\), both Fibonacci
factors are positive. Theorem~\ref{thm:signed-gap}
therefore proves the parity rule.

For the three families choose \(c=1\).
Using the elementary identities
\(L_{n+1}-L_n=L_{n-1}\),
\(L_{n+2}-L_n=L_{n+1}\), and
\(L_{n+4}-L_n=5F_{n+2}\), take
\[
(h,a)=(1,3r+1),\qquad
(2,4r+1),\qquad
(4,3r-2),
\]
respectively (the last has \(r\ge2\)).
The denominators \(D(1,1+h)\) are \(2,3,10\).
Divisibility follows from the period-three
Lucas/Fibonacci parities and the Lucas residues
\(L_{4r+2}\equiv0\pmod3\).
Theorem~\ref{thm:global} proves that every
fixed \(S_q\) has only finitely many collisions.
\end{proof}

\section*{A family with globally unique differences}
\begin{theorem}\label{thm:unique}
Fix $r\ge1$ and set $q=5F_{2r}$. Then
\[
L_i-qL_j>0\quad\Longleftrightarrow\quad
(i,j)=(2r+2,1)\quad\text{or}\quad
i\ge2r+3,\;1\le j\le i-2r-2.
\]
Every positive integer in $S_q$ has precisely one representation.
\end{theorem}
\begin{proof}
The recurrence yields
\begin{equation}\label{eq:q}
q=5F_{2r}=L_{2r+2}-L_{2r-2},
\quad q-L_{2r+1}=L_{2r-1}>0,
\quad L_{2r+2}-q=L_{2r-2}>0.
\end{equation}
Thus $j=1$ first becomes admissible at $i=2r+2$.
For $j\ge2$ put
$A_j=L_{j+2r+2}-qL_j$ and
$B_j=qL_j-L_{j+2r+1}$.
Both obey the Lucas recurrence in $j$. Checking two initial values
in $r$ and using the recurrence in steps of two gives
\[
A_2=L_{2r+4}-3q=L_{2r-4}>0,
\quad A_3=A_2+A_1>0,\quad A_1=L_{2r+3}-q>0,
\]
with $L_{-2}=3$ when $r=1$.
For $B_j$, the first two terms of the sequences $B_2(r)$
and $B_3(r)$ are
\[
(B_2(1),B_2(2))=(4,16),\qquad
(B_3(1),B_3(2))=(2,13).
\]
Each satisfies $b_{r+2}=3b_{r+1}-b_r$ as a function of $r$.
An increasing positive initial pair remains increasing and positive.
Therefore $B_2,B_3>0$ for all $r\ge1$; the recurrence in $j$
then gives $A_j,B_j>0$ for all $j\ge2$.
Hence $L_{j+2r+1}<qL_j<L_{j+2r+2}$ for $j\ge2$,
establishing the stated index region.

Let $\mathcal B_i$ be the positive values with first index $i$.
They are distinct within each block, and
$\max\mathcal B_i=L_i-q$ for $i\ge2r+2$.
Using \eqref{eq:identities} with $b=i-2r+2$, the next block
of the same residue class satisfies
\begin{align}\label{eq:gap}
\min\mathcal B_{i+4}-\max\mathcal B_i
&=L_{i+4}-L_i-qL_{i-2r+2}+q\\
&=q+5F_{i-4r+2}>0.\nonumber
\end{align}
For a nonnegative Fibonacci index the inequality is immediate.
For a negative index, $i\ge2r+2$ implies
$0<|i-4r+2|\le2r-4<2r$, so
$|F_{i-4r+2}|<F_{2r}$.
Moreover $A_{j+4}>A_j$ follows from the strictly positive
Lucas recurrence; thus blocks separated by any positive multiple
of four are strictly ordered.

Finally, $(L_1,L_2,L_3,L_4)\bmod5=(1,3,4,2)$
is the repeating residue cycle. Since $5\mid q$,
two equal differences must have first indices congruent modulo four.
They cannot belong to different such blocks, by \eqref{eq:gap},
and the same block has distinct values.
\end{proof}



\section*{The four OEIS sequences at coefficient five}
The preceding classification applies in particular to \(q=5\).
For this section write \(S=S_5\), \(T=T_5\),
and \(P=P_5\). Their increasing positive values
give OEIS A191981, A191982, and A191980,
respectively; the positive values of A191983 form
\(S\cup T\). The published A191981 and A191983
also begin with a zero sentinel. Because
\(L_i\bmod5\) cycles through \(1,3,4,2\),
the equation \(L_i=5L_j\) is impossible:
the zero is not a Lucas difference.

Theorem~\ref{thm:unique} with \(r=1\) gives
the complete index region and uniqueness in \(S\):
\[
(i,j)=(4,1)\quad\text{or}\quad
i\ge5,\ 1\le j\le i-4.
\]
For \(i\ge5\) its ordered block satisfies
\begin{equation}\label{eq:oeis-block}
3L_{i-5}=L_i-5L_{i-4}
\le L_i-5L_j\le L_i-5.
\end{equation}
Write \(H(y)=\#\{j\ge1:L_j\le y\}\) and
\([x]_+=\max(x,0)\). The exact conclusions
for the four sequences are as follows.

\subsection*{Exact positive-difference count}
\begin{theorem}[Three-block exact counting]\label{oeis:thm:exact}
For every integer $x\ge11$, let $m$ be the unique integer satisfying
$L_m\le x<L_{m+1}$. Then
\begin{equation}\label{oeis:eq:exact}
\boxed{\displaystyle A(x)=1+\frac{(m-4)(m-3)}2+
\sum_{h=1}^{3}\left[m+h-4-
H\!\left(\left\lfloor\frac{L_{m+h}-x-1}{5}\right\rfloor\right)
\right]_+.}
\end{equation}
For $0\le x<11$, $A(x)$ counts the members of $\{2,3,6,9\}$ not exceeding $x$.
\end{theorem}
\begin{proof}
All admissible differences with first index $i\le m$ are smaller than
$L_i\le x$. For $i\ge m+4$, \eqref{eq:oeis-block} gives
\[
L_i-5L_j\ge3L_{i-5}\ge3L_{m-1}>L_{m+1}>x.
\]
Indeed, $3L_{m-1}-L_{m+1}=L_{m-1}-L_{m-2}>0$ because $m\ge5$.
The only remaining first indices are $i=m+1,m+2,m+3$. In one such block,
$L_i-5L_j\le x$ is equivalent to
$L_j>\lfloor(L_i-x-1)/5\rfloor$.
There are $i-4$ admissible indices $j$, so removing the ineligible ones gives
the positive-part expression in \eqref{oeis:eq:exact}.
\end{proof}

\subsection*{Complete intersection of the two differences}
\begin{theorem}[Complete intersection classification]\label{oeis:thm:intersection}
With $E=\{2,6,9,12,13,14,24,32,61\}$, one has the disjoint union
\begin{equation}\label{oeis:eq:intersection}
\boxed{\quad S\cap T=E\ \dot\cup\
\{3L_n:n\ge6\}\ \dot\cup\ \{8L_n:n\ge6\}.\quad}
\end{equation}
In particular, the only common values outside the two Lucas-multiple
families are the nine elements of $E$.
\end{theorem}
\begin{proof}
Suppose $L_i-5L_j=5L_k-L_\ell>0$.
The Lucas residues modulo $5$ repeat as $(1,3,4,2)$, so
$L_i+L_\ell\equiv0\pmod5$ forces
$\ell\equiv i+2\pmod4$.
By Theorem~\ref{thm:unique}, $i\ge4$ and $j\le i-4$, except
that $(i,j)=(4,1)$ is the first permitted pair. Furthermore,
$5L_k>L_\ell$ implies $\ell\le k+3$ for every $k\ge1$.
For $k\ge2$, the least positive difference in its block is
$5L_k-L_{k+3}=2L_{k-2}$.

Assume first $i\ge16$. The lower bound
$L_i-5L_j\ge3L_{i-5}>5L_{i-7}$ excludes $k\le i-7$.
The bound $5L_k-L_\ell\ge2L_{k-2}\ge2L_{i-1}>L_i$
excludes $k\ge i+1$. If $k=i$, the largest permissible
$\ell$ is $i+2$ and $5L_i-L_{i+2}>L_i$.
If $k=i-2$, the largest permissible $\ell$ is $i-2$ and
$4L_{i-2}>L_i$. Thus $k=i-d$ with $d\in\{1,3,4,5,6\}$.
For $d=1$, only $\ell=i+2$ can give a value below $L_i$;
all smaller permissible $\ell$ give $5L_{i-1}-L_\ell>L_i$.

Put $R=L_i+L_\ell-5L_{i-d}=5L_j$.
The following table gives strict open intervals for $R$;
each endpoint is a multiple of a Lucas number.
\[
\begin{array}{c|c|l}
d&\ell&\text{location of }R\\ \hline
1&i+2&(5L_{i-5},\,5L_{i-4})\\
3&i-2&(5L_{i-7},\,5L_{i-6})\\
3&\ell\le i-6&R<0\\
4&i-2&(5L_{i-5},\,5L_{i-4})\\
4&i-6&(5L_{i-6},\,5L_{i-5})\\
4&i-10&R=5L_{i-6}\\
4&\ell\le i-14&(5L_{i-7},\,5L_{i-6})\\
5&i-2&R>5L_{i-4}\\
5&\ell\le i-6&(5L_{i-5},\,5L_{i-4})\\
6&i-6&R>5L_{i-4}\\
6&i-10&R=5L_{i-4}\\
6&\ell\le i-14&(5L_{i-5},\,5L_{i-4})
\end{array}
\]
All table inequalities hold for $i\ge16$.
To verify them, bound each variable $\ell$ between its
indicated upper endpoint and $L_\ell>0$.
The resulting fixed-shift inequalities are differences of
Lucas-linear expressions and hence satisfy
$u_{i+2}=u_{i+1}+u_i$. In each case the two initial values
$i=16,17$ are strictly positive, as checked by direct substitution;
the recurrence preserves positivity. The two exact-equality
rows follow from the same recurrence and two initial identities.
Every interval either misses all admissible $5L_j$ or contains
no such value, while $R>5L_{i-4}$ is impossible since $j\le i-4$.
The only surviving cases are
\[
(i,j,k,\ell)=(i,i-6,i-4,i-10),\qquad
(i,i-4,i-6,i-10)\qquad(i\ge16).
\]

For the remaining $4\le i\le15$, the finite possibilities
$1\le k\le i-1$, $1\le\ell\le k+3$, and
$j\le i-4$ (with $(4,1)$ included) give:
\[
\begin{array}{c|ccccccc}
i&4&5&6&7&8&9&10\\ \hline
S\cap T\text{ in the }i\text{-block}
&2&6&13&9,14,24&12,32&61&\varnothing
\end{array}
\]
For $11\le i\le15$ the only two values are
$3L_{i-5}$ and $8L_{i-5}$.
For all $n\ge6$ the Lucas recurrence gives
\begin{align*}
3L_n&=L_{n+5}-5L_{n+1}=5L_{n-1}-L_{n-5},\\
8L_n&=L_{n+5}-5L_{n-1}=5L_{n+1}-L_{n-5}.
\end{align*}
Thus both displayed families belong to $S\cap T$.
They are disjoint from one another by the uniqueness theorem
for $S$, and neither family contains one of the nine exceptional
values (the smallest values are $3L_6=54$ and $8L_6=144$). This proves \eqref{oeis:eq:intersection}.
\end{proof}

\subsection*{All repeated reverse differences and sums}
\begin{theorem}[Five and only five collisions]\label{oeis:thm:duplicates}
The map
\[
(k,\ell)\longmapsto 5L_k-L_\ell>0,\qquad k,\ell\ge1,
\]
is injective except for exactly five pairs of distinct arguments:
\[
\begin{array}{c|c|c|c}
k&p&(\ell,q)&5L_k-L_\ell=5L_p-L_q\\ \hline
1&2&(1,5)&4\\
1&3&(2,6)&2\\
2&5&(4,8)&8\\
2&6&(1,9)&14\\
3&4&(2,6)&17
\end{array}
\]
In particular, the set of multiply represented positive values in
A191982 is precisely $D=\{2,4,8,14,17\}$; each has multiplicity two.
\end{theorem}
\begin{proof}
For $k=1$ the positive pairs have $1\le\ell\le3$.
For every $k\ge2$, positivity is equivalent to
$1\le\ell\le k+3$, and the values in the $k$-block lie in
\[
[2L_{k-2},\,5L_k-1],
\qquad 5L_k-L_{k+3}=2L_{k-2}.
\]
Suppose $5L_k-L_\ell=5L_p-L_q>0$ with $k<p$.
Then
\begin{equation}\label{oeis:eq:collision}
L_q-L_\ell=5(L_p-L_k)>0,\qquad q\equiv\ell\pmod4.
\end{equation}
The congruence follows from the four distinct Lucas residues modulo $5$.
Hence $q\ge\ell+4$.

If $k\ge3$ and $p\ge k+4$, then
\[
5L_p-L_q\ge2L_{p-2}\ge2L_{k+2}>5L_k-1,
\]
where $2L_{k+2}-5L_k=L_{k-1}-L_{k-2}>0$.
This exceeds every value in the $k$-block.
For $k=1$, $p\ge5$ is excluded by
$2L_3=8>5L_1-1=4$.
For $k=2$, $p\ge7$ is excluded by
$2L_5=22>5L_2-1=14$; the boundary $p=6$
gives precisely $14$.

It remains to consider $p=k+d$ for $d\in\{1,2,3\}$.
For $k\ge4$ the possibilities in \eqref{oeis:eq:collision}
reduce as follows.

If $d=1$, then $5(L_p-L_k)=5L_{k-1}$,
so $L_{k+2}<5L_{k-1}<L_{k+4}$ and
$q\in\{k+3,k+4\}$.
For $q=k+3$ one has
$L_\ell=3L_{k-2}>L_{k-1}$, contradicting
$\ell\le q-4=k-1$.
For $q=k+4$ one has
$L_\ell=L_{k+2}+3L_{k-2}>L_k$,
contradicting $\ell\le k$.

If $d=2$, then $5(L_p-L_k)=5L_{k+1}$.
Since $L_{k+4}<5L_{k+1}$ and $q\le p+3$,
necessarily $q=k+5$.
It follows that
$L_\ell=3L_k>L_{k+1}$, but $\ell\le k+1$.

If $d=3$, then $5(L_p-L_k)=10L_{k+1}$.
As $L_{k+5}<10L_{k+1}$ and $q\le p+3$,
one must have $q=k+6$. Thus
\[
L_\ell=L_{k+6}-10L_{k+1}
=5L_k-2L_{k+1}=:R_k.
\]
For $k\ge3$,
\[
L_{k-2}<R_k<L_{k+2}.
\]
But \eqref{oeis:eq:collision} forces
$\ell\equiv k+2\pmod4$ and $\ell\le k+2$;
either $\ell=k+2$ or $\ell\le k-2$, a contradiction.

Finally, the cases $k=1,2,3$ can be checked using
\eqref{oeis:eq:collision} and $q\le p+3$:
$k=1$ yields $(p,\ell,q)=(2,1,5),(3,2,6)$;
$k=2$ yields $(5,4,8)$ in addition to the
boundary $(6,1,9)$; $k=3$ yields $(4,2,6)$.
Their values are the five displayed in the table.
Since the list exhausts all colliding argument pairs,
no value has three representations.
\end{proof}

\begin{theorem}[The five sum collisions]\label{oeis:thm:sums}
Exactly five values in $P$ admit two different index-pair
representations:
\[
\begin{array}{c|l}
s&\text{the two representations of }s\\ \hline
16&5L_1+L_5=5L_2+L_1\\
23&5L_1+L_6=5L_3+L_2\\
38&5L_3+L_6=5L_4+L_2\\
62&5L_2+L_8=5L_5+L_4\\
91&5L_2+L_9=5L_6+L_1
\end{array}
\]
All other values have precisely one representation; none has three.
\end{theorem}
\begin{proof}
Suppose $5L_i+L_j=5L_k+L_\ell$ with $(i,j)\ne(k,\ell)$.
Since $L_n$ increases strictly for $n\ge1$, we may take
$i<k$ and $j>\ell$. Rearranging gives
\begin{equation}\label{oeis:eq:sum-switch}
5L_i-L_\ell=5L_k-L_j=:t.
\end{equation}
The four Lucas residues $1,3,4,2$ modulo $5$ imply
$t\ne0$. If $t<0$, then
$-t=L_\ell-5L_i=L_j-5L_k>0$
has two distinct representations in $S$, contradicting
Theorem~\ref{thm:unique}. Consequently $t>0$.
Equation \eqref{oeis:eq:sum-switch} is therefore a collision
in $T$, and Theorem~\ref{oeis:thm:duplicates} restricts it to
the five pairs listed there.
Conversely, swapping the subtracted Lucas indices in
each of those five equalities yields exactly the five
sum identities displayed above. Each displayed sum is distinct,
and each of the five collisions in $T$ is between precisely
two index pairs; hence there are no further collisions and
no threefold representations in $P$.
\end{proof}

\subsection*{Complete intersection of sums and reverse differences}
\begin{theorem}[Complete sum--difference intersection]\label{oeis:thm:sum-full-overlap}
The following union is disjoint:
\begin{equation}\label{oeis:eq:sum-full-overlap}
\boxed{P\cap T=E\ \dot\cup\
 \{5L_n+L_{n+7}:n\ge4\}\ \dot\cup\
 \{5L_n+L_{n-3}:n\ge4\}.}
\end{equation}
Indeed, every equality $5L_a+L_b=5L_c-L_d>0$
with $c\ge8$ has $a=c-4$ and
$\{b,d\}=\{c+3,c-7\}$.
\end{theorem}
\begin{proof}
The equality implies $a<c$, since $L_b+L_d>0$, and
\begin{equation}\label{oeis:eq:sum-overlap-condition}
L_b+L_d=5(L_c-L_a),\qquad b\equiv d+2\pmod4.
\end{equation}
The residue condition follows from the four distinct Lucas
residues modulo $5$. Positivity gives $d\le c+3$; also
$L_b<5L_c<L_{c+4}$ gives $b\le c+3$.

Assume $c\ge8$, and set $m=\max(b,d)$ and $s=\min(b,d)$.
The residue condition forces $m-s\equiv2\pmod4$,
so $s\le m-2$. If $m=c+3$, then $s=c+1$ would
make $L_m+L_s>5L_c$, hence $s\le c-3$.
If $m=c+2$, then $s\le c$; if $m\le c+1$, then
$s\le c-1$. The Lucas recurrence yields
\[
L_m+L_s\le L_{c+3}+L_{c-3},\qquad
5L_c-L_{c+3}-L_{c-3}>5L_{c-5}.
\]
Thus $L_a>L_{c-5}$, hence $a\ge c-4$.
Write $h=c-a\in\{1,2,3,4\}$ and
$R_h=5(L_c-L_{c-h})=L_m+L_s$.

For $h=1$, the bounds
$2L_{c-1}<R_1<L_{c+2}$ force $m=c$ or $c+1$.
When $m=c$, $R_1>L_c+L_{c-2}$ contradicts $s\le c-2$.
When $m=c+1$,
\[
L_{c-5}<R_1-L_{c+1}<L_{c-1}.
\]
But then $s$ must belong to $\{c-1,c-5,c-9,\ldots\}$,
which is impossible by strict monotonicity.
For $h=2$, the bounds $2L_c<R_2<L_{c+3}$
force $m=c+1$ or $c+2$. The former is excluded by
$R_2>L_{c+1}+L_{c-1}$; in the latter case
\[
L_{c-4}<R_2-L_{c+2}<L_c,
\]
whereas $s\in\{c,c-4,c-8,\ldots\}$.
For $h=3$, $2L_{c+1}<R_3<L_{c+3}$ forces
$m=c+2$, but $R_3>L_{c+2}+L_c$ contradicts $s\le c$.
Finally, for $h=4$, $R_4>L_{c+2}+L_c$ forces
$m=c+3$; the identity
\[
5(L_c-L_{c-4})=L_{c+3}+L_{c-7}
\]
gives $s=c-7$. This gives precisely the two families
in \eqref{oeis:eq:sum-full-overlap}.

All strict inequalities used above are fixed-shift linear
combinations of Lucas numbers. Their differences satisfy
$u_{c+2}=u_{c+1}+u_c$; each is positive at $c=8,9$,
so all hold for $c\ge8$. The companion verifier records
the 18 positive initial comparisons and the exact identity.

For $c\le7$, the ranges $1\le a<c$ and
$1\le b,d\le c+3$ are finite. Exhaustion gives the
following distinct values, grouped by $c$:
\[
\begin{array}{c|l}
c &5L_a+L_b=5L_c-L_d\ \text{(distinct values)}\\ \hline
1&\varnothing\\
2&8,12\\
3&9,16,19\\
4&6,24,31,34\\
5&8,26,44,52\\
6&\varnothing\\
7&22,138
\end{array}
\]
These give exactly $E$. The family with $L_{n+7}$
begins at $234>\max E$; the other begins with
$36,58,94,152,\ldots$, none of which belongs to $E$.
The two families are disjoint by Theorem~\ref{oeis:thm:sums}:
if one value belonged to both, it would have two different
sum representations, yet every value of the first family
exceeds $91$. The Lucas identity
\[
5(L_{n+4}-L_n)=L_{n+7}+L_{n-3}\qquad(n\ge4)
\]
holds by the recurrence and its first two instances.
It gives both families in \(P\cap T\);
the second difference is positive since
\(5L_{n+4}-L_{n+7}=2L_{n+2}>0\).
\end{proof}

\subsection*{Exact counting and unranking}
Let
\[
A(x)=\#(S\cap[1,x]),\quad
B(x)=\#(T\cap[1,x]),\quad
C(x)=\#(P\cap[1,x]).
\]
The first count is \eqref{oeis:eq:exact} for \(x\ge11\);
for \(0\le x<11\), count the elements of
\(\{2,3,6,9\}\) at most \(x\).
Let
\[
D_T=\{2,4,8,14,17\},\qquad
D_P=\{16,23,38,62,91\},
\]
and put \(b_1=3\), \(b_k=k+3\) for \(k\ge2\),
\(K(x)=\max\{3,H(\lfloor x/2\rfloor)+2\}\).
The two remaining count formulas are
\begin{align}
B(x)&=\sum_{k=1}^{K(x)}
[b_k-H(5L_k-x-1)]_+
-\#\{d\in D_T:d\le x\},\label{eq:oeis-Tcount}\\
C(x)&=\sum_{i=1}^{H(\lfloor(x-1)/5\rfloor)}
H(x-5L_i)
-\#\{d\in D_P:d\le x\}.\label{eq:oeis-Pcount}
\end{align}
For the two intersections let
\(I(x)=\#(S\cap T\cap[1,x])\) and
\(J(x)=\#(P\cap T\cap[1,x])\). Write
\[
E_{ST}=\{2,6,9,12,13,14,24,32,61\},\qquad
E_{PT}=\{6,8,9,12,16,19,22,24,26,31,34,44,52,138\},
\]
and \(u_n=5L_n+L_{n+7}\),
\(v_n=5L_n+L_{n-3}\) for \(n\ge4\).
Then the two intersection theorems yield
\begin{align}
I(x)&=\#\{e\in E_{ST}:e\le x\}
+[H(\lfloor x/3\rfloor)-5]_+
+[H(\lfloor x/8\rfloor)-5]_+,\label{eq:oeis-Icount}\\
J(x)&=\#\{e\in E_{PT}:e\le x\}
+\#\{n\ge4:u_n\le x\}
+\#\{n\ge4:v_n\le x\}.
\label{eq:oeis-Jcount}
\end{align}
Finally, for \(x\ge0\), the number of entries
of A191983 not exceeding \(x\) is
\begin{equation}\label{eq:oeis-Ucount}
\boxed{\,U(x)=1+A(x)+B(x)-I(x).\,}
\end{equation}
Every sum in these formulas is finite.
They count pairs within their permitted Lucas
index regions, then subtract the five duplicated
values in \(T\) or \(P\), respectively.
The intersection counts follow from the disjoint
unions proved above, and \(U\) follows from
inclusion--exclusion. All four ordered sequences
and both intersections therefore admit exact
unranking by integer binary search.

\paragraph{OEIS references.}
The four records are
\url{https://oeis.org/A191980},
\url{https://oeis.org/A191981},
\url{https://oeis.org/A191982}, and
\url{https://oeis.org/A191983}.

\end{document}